\section*{Executive summary}\label{sec:summary}
\addcontentsline{toc}{section}{Executive summary}
\markboth{Executive summary}{Executive summary}

\subsection*{The result}

\evobs{} The absence check ties its mismatched-evidence control; prediction of hidden knowledge is untested. The
proof of concept (PoC) built a chain from public web text to ranked questions for experts. Its
central step, the absence check (\emph{closure} in the code), asks a judge model whether retrieved
ledger sentences already state an element that a lens flagged as missing. In a control, each
candidate's additional other-source sentences were swapped for those of the topically nearest
candidate, while the seed's own sentences stayed in both arms. The judge's closure rate (counting
\emph{partial} as closed) did not change: 0.889, 0.961 and 0.958 in the three ledgers. Under this
control, the judge did not tell a candidate's own cross-source evidence from the donor's, so it is not
yet an informative absence check; the test does not show why. The control is built into \code{gapmap/checks.py}, added after an
adversarial AI review found the weakness, and it flags the failure on every replay. No gold standard
was acquired, so the PoC neither shows nor refutes that the gap map predicts what experts know. A later
exploratory spike with a low-cost decision model (Jev) on the committed N3 output gave no per-record signal under its pre-set form test (the score tracked text form). After its first pre-set judge rule failed, it agreed with a second reference, a stronger model, better than the current closure judge did; it detects no hidden knowledge (\cref{sec:jev}).

\subsection*{Recommendation: run the closure study before any grant submission}

\evfut{} The cheapest decisive test needs no grant and no participants. In this Phase~0 study (V1),
blind coders who are not the owner (two preferred) label each unit as closed, partly closed or open. A unit is one
lens firing with its evidence, in one arm of the mismatched-evidence control. The sample is the 107 lens firings from the two admissible ledgers (105 with a donor, in both arms: 210 pools; the firing is the unit of inference), plus about 100 units from one new
domain, so that no judge is scored only in-sample. V1 reports three things:
\begin{itemize}[nosep]
  \item agreement between the two human coders (weighted Cohen's $\kappa$; the gate uses closed
    versus not closed);
  \item corpus-level retrieval recall: is the element stated anywhere in the ledger's sources, not
    only in the sentences the retriever returned;
  \item whether any judge in a candidate panel beats the mismatched-evidence control.
\end{itemize}
The cost is tens of coder hours plus local compute, on the order of a few thousand euros. The gate
MS1 reads CONTINUE only if human $\kappa \geq 0.70$, retrieval recall meets its pre-registered floor,
and at least one judge beats the mismatched-evidence control by a pre-registered margin. Otherwise the automated
absence check stops, and the work pivots to reconstruction with provenance as an audit tool, or to
interviewing without a gap map. The V1 numbers should be on page~1 of any grant proposal; this report
is the design for them (\cref{sec:validation}).

\subsection*{The problem and the idea}

Experts leave out the steps that have become automatic for them. In one small study, three
surgeons teaching a procedure omitted on average 73\% of the decision steps in a task list built by
cognitive task analysis (CTA) \parencite{sullivan2014use}; no general omission rate is known. Text
has the same silence: writers do not state what is obvious to them \parencite{gordon2013reporting}.

We ask AI for two narrow things. First, \emph{reconstruct} what public sources state about a task,
as a ledger of labeled claims whose evidence items are exact spans of saved sources. Second,
\emph{predict where that reconstruction is thin}, using constructs from expertise research as
methodology lenses. A lens asks: \enquote{the sources name this cue, rival cause, rule or check; do
they ever state its content?} What survives the absence check becomes a \emph{gap candidate}: a
hypothesis labeled \emph{inferred}, with a template question for an expert. What experts would add beyond the ledger is
the Human Knowledge Residual, a construct nobody has measured yet.

\subsection*{What the PoC built}

\evobs{} The PoC used zero participants and public web text only, in two domains: intermittent
faults in programmable logic controllers (PLC) and data protection impact assessments under the EU General Data Protection Regulation
(GDPR), at
\$0.50--\$1.00 per run.

\begingroup
  \footnotesize\noindent
  \begin{tabularx}{\linewidth}{@{}>{\bfseries}p{2.1cm}X@{}}
    \toprule
    Built & Evidence model with computed labels and a frozen schema (N1); reconstruction pipeline
      with verbatim spans and cross-family verification (N2); gap map with seven lenses, absence
      check, triage, a breadth ranking (not expected information gain), template questions and
      built-in checks (N3). \\
    Shown & The chain runs end to end on real sources and replays byte for byte. All 909 stored
      evidence items re-slice exactly from their snapshots; 260 of 1,169 claims have no locatable span and 158 more have a span the verifier did not accept; all 418 are labeled synthetic. On the two admissible ledgers the lenses fired 107 times and
      left 25 gap candidates after merging (PLC 24, GDPR v1 1); a third, defect-compromised ledger
      adds 48 firings and 6 candidates. \\
    Not shown & That any gap candidate is real expert knowledge; that the absence check works; that
      gating resists planted falsehoods (N2's gated arms were never scored, so its registered verdict
      is inconclusive); that the ranking is more than source density; that anything generalizes
      beyond two domains. \\
    \bottomrule
  \end{tabularx}
\endgroup

\paragraph{What is new, and what is not.} \evint{} Every step has precedent. The closest work checks
GDPR privacy policies for typed, regulation-required content \parencite{torre2020ai}, finds
underspecified phrases in instructional text \parencite{anthonio2022clarifying}, and mines design
rationale from issue trackers \parencite{zhao2024drminer}. What we did not find, in a search to
September 2026, is a method that chains a provenance-typed ledger, expertise constructs as absence
detectors, a knowledge type and elicitation channel per candidate, and search failures routed back
to search. That chain is a proof of concept, not a validated method (\cref{sec:approach-closest}).

\subsection*{The program after V1}

\evfut{} Each tranche of money is released only by a gate reading CONTINUE. An external gate reader,
an independent statistician, reads the three deciding gates (MS1, MS3, MS5) against their
pre-registration; any exception by the owner, at any gate, needs the reader's sign-off.
\begin{itemize}[nosep]
  \item \textbf{Phase~1, a lean methods grant} (about EUR 0.3--0.5 million, 12--15 months), only on
    MS1 CONTINUE. It builds an area score and freezes the configuration before any gold is acquired.
    It then runs V2: the gap map is scored against several published CTA task lists, pooled, on
    dated corpora. The gate MS3 only screens money for Phase~2; it does not decide H2, the central
    hypothesis that the gap map predicts where experts add knowledge. The score is the gain in
    area under the receiver operating characteristic (ROC) curve ($\Delta$AUROC) over the strongest baseline. MS3 releases Phase~2 only
    if its point estimate is at least $\delta$ (proposed 0.05), its lower 90\% bound is above
    $-\delta/2$, and the map adds value beyond source density and area size. Otherwise H2 funding stops and the
    separately costed interview-only pivot follows.
  \item \textbf{Phase~2, the full program}, only on MS3 CONTINUE. Blind prospective elicitation
    (V5), with a provisional smallest gain of interest $\delta = 0.05$ and a size set after its
    pilot, is the prospective test of H2. H2 stops if the upper 90\% bound of $\Delta$AUROC is below $\delta$; any reading
    short of CONTINUE does not support H2 and ends its funding. A later study tests H3: prioritized
    questions recover at least 1.25 times as many validated items per expert hour; a null counts
    against it. Organizational and education pilots follow, each with its own gate.
\end{itemize}
A study of developer departures in open-source projects tests a separate organizational hypothesis
(H-OSS) and can neither stop nor license H2. Budgets: \cref{sec:resources}. \textbf{Team status:} one operator
built the PoC with AI assistance; a host institution, principal investigator, statistician, CTA
collaborator and letters must be named before submission (\cref{sec:team}).
